% Worked example, version A: six panels in reading order, two rows (full text width).
\documentclass[border=3pt]{standalone}
% Shared preamble of the worked-example figures (case_a/b/c.tex); data from ../generated/case_data.tex
\usepackage{libertinus}
\renewcommand{\familydefault}{\sfdefault}
\usepackage{tikz}
\usepackage{amsmath}
\usepackage{booktabs}
\usetikzlibrary{arrows.meta,positioning,calc,fit,backgrounds,decorations.pathreplacing}
\makeatletter
\newcommand{\C}[1]{\@ifundefined{case@#1}{\textcolor{red}{[#1?]}}{\csname case@#1\endcsname}}
\makeatother
% word-class palette (decision_shares.png / playground), one slot per category
\definecolor{catF}{HTML}{2A78D6}  % function words
\definecolor{catV}{HTML}{EB6834}  % verbs
\definecolor{catN}{HTML}{1BAF7A}  % nouns and realia
\definecolor{catA}{HTML}{EDA100}  % agents and deities
\definecolor{catX}{HTML}{E87BA4}  % foreign words
\definecolor{catO}{HTML}{9A9A94}  % other / unclassified
\definecolor{shade}{HTML}{EB6834}
\definecolor{ink}{HTML}{2B2B2B}
\definecolor{soft}{HTML}{8A8A8A}
\definecolor{topc}{HTML}{1F7A6B}
\definecolor{secc}{HTML}{5B4B8A}
% a word tinted by its class; dashed frame when the class is low-confidence
\tikzset{wc-low/.style={densely dashed}, wc-ok/.style={}, reset/.style={text width={}, minimum height=0pt, minimum width=0pt, align=center, draw=none, fill=none, solid, font=}}
\newcommand{\wc}[3]{\tikz[baseline=(w.base)]\node[reset, inner sep=0.8pt, rounded corners=1pt, fill=#1!22,
  draw=#1!70, line width=0.4pt, wc-#2](w){\strut #3};}
% a word shaded by its share of the decision (0-100)
\newcommand{\ws}[2]{\tikz[baseline=(w.base)]\node[reset, inner sep=0.8pt, rounded corners=1pt, fill=shade!#1](w){\strut #2};}
\newcommand{\broken}[1]{\textcolor{soft}{#1}}
\newcommand{\tline}[2]{\noindent\hangindent=5.7mm\hangafter=1\makebox[4.5mm][r]{\textcolor{soft}{#1}}\hspace{1.2mm}#2\par}
\newcommand{\wf}[1]{\textbf{#1}}          % a word form
\newcommand{\sg}[1]{\textit{#1}}          % a sign sequence
\newcommand{\stepno}[1]{\tikz[baseline=-0.6ex]\node[reset, circle, fill=ink, text=white, inner sep=0.6pt, font=\scriptsize\bfseries]{#1};}
\newcommand{\legend}{%
  \wc{catF}{ok}{function} \wc{catV}{ok}{verb} \wc{catN}{ok}{noun, realia} \wc{catA}{ok}{agent, deity}
  \wc{catX}{ok}{foreign} \wc{catO}{low}{unclear}}
% generated by make_case_figures.py; do not edit
\expandafter\def\csname case@txtid\endcsname{KBo 41.22}
\expandafter\def\csname case@filed\endcsname{470}
\expandafter\def\csname case@words\endcsname{33}
\expandafter\def\csname case@top\endcsname{Luwian rituals (Puriyanni, Kuwattalla)}
\expandafter\def\csname case@topshort\endcsname{Luwian rituals}
\expandafter\def\csname case@second\endcsname{Kizzuwatna rituals}
\expandafter\def\csname case@secondshort\endcsname{Kizzuwatna rituals}
\expandafter\def\csname case@s1\endcsname{0.39}
\expandafter\def\csname case@s2\endcsname{0.25}
\expandafter\def\csname case@margin\endcsname{0.14}
\expandafter\def\csname case@tau\endcsname{0.018}
\expandafter\def\csname case@nfeat\endcsname{177}
\expandafter\def\csname case@dim\endcsname{25{,}372}
\expandafter\def\csname case@ntop\endcsname{10}
\expandafter\def\csname case@nsecond\endcsname{165}
\expandafter\def\csname case@common\endcsname{−21\%}
\expandafter\def\csname case@idf.rare\endcsname{6.8}
\expandafter\def\csname case@idf.common\endcsname{1.2}
\expandafter\def\csname case@nn1\endcsname{KBo 9.125+}
\expandafter\def\csname case@nn1cth\endcsname{409}
\expandafter\def\csname case@nn.409\endcsname{5}
\def\CaseText{\tline{3′}{\wc{catJ}{ok}{i-da-a-lu-un} \wc{catO}{low}{UD-an} \wc{catO}{low}{ma-}}\tline{4′}{\wc{catA}{ok}{DINGIR\textsuperscript{MEŠ}-aš} \wc{catT}{ok}{kar-pí-in} \wc{catO}{low}{pa-an-ga-u-wa-aš} \wc{catO}{low}{EME-an}}\tline{5′}{\wc{catV}{ok}{ša-me-ed-du-wa-ra-at-kán} \broken{ḫi-x}}\tline{6′}{\wc{catF}{ok}{na-aš-ta} \wc{catA}{ok}{EN} \wc{catN}{ok}{SÍSKUR} \wc{catT}{ok}{ŠUM-ŠU}}\tline{7′}{\wc{catF}{ok}{le-e} \wc{catO}{low}{ku-wa-pí-ik-ki} \broken{a-x}}\tline{8′}{\wc{catT}{ok}{te-eš-ḫa-an} \wc{catO}{low}{ma-ni-ku-wa-an(-)}}\tline{9′}{\wc{catF}{ok}{nu-uš-ši} \wc{catN}{ok}{\textsuperscript{GI}pát-tar} \wc{catF}{ok}{pa-ra-a}}\tline{10′}{\wc{catA}{ok}{\textsuperscript{MUNUS}ŠU.GI-ma} \wc{catV}{ok}{me-ma-i} \broken{x}}\tline{11′}{\wc{catA}{ok}{DINGIR\textsuperscript{MEŠ}-aš} \wc{catT}{ok}{kar-pí-in} \wc{catO}{low}{pa-an-ga-u-wa-aš}}\tline{12′}{\wc{catA}{ok}{\textsuperscript{LÚ}A.ZU-ma-aš-ši-iš-ša-an}}\tline{13′}{\wc{catF}{ok}{na-aš} \wc{catF}{ok}{ša-ra-a} \wc{catO}{low}{ti-}}\tline{14′}{\wc{catF}{ok}{na-aš-ta} \wc{catB}{ok}{GU₄-un}}\tline{15′}{\wc{catF}{ok}{a-ra-aḫ-za-an-da}}}
\def\CaseShaded{\tline{3′}{\ws{27}{i-da-a-lu-un} \ws{2}{UD-an} \ws{0}{ma-}}\tline{4′}{\ws{2}{DINGIR\textsuperscript{MEŠ}-aš} \ws{39}{kar-pí-in} \ws{70}{pa-an-ga-u-wa-aš} \ws{93}{EME-an}}\tline{5′}{\ws{0}{ša-me-ed-du-wa-ra-at-kán} \broken{ḫi-x}}\tline{6′}{\ws{0}{na-aš-ta} \ws{0}{EN} \ws{0}{SÍSKUR} \ws{14}{ŠUM-ŠU}}\tline{7′}{\ws{0}{le-e} \ws{15}{ku-wa-pí-ik-ki} \broken{a-x}}\tline{8′}{\ws{0}{te-eš-ḫa-an} \ws{6}{ma-ni-ku-wa-an(-)}}\tline{9′}{\ws{0}{nu-uš-ši} \ws{0}{\textsuperscript{GI}pát-tar} \ws{2}{pa-ra-a}}\tline{10′}{\ws{27}{\textsuperscript{MUNUS}ŠU.GI-ma} \ws{0}{me-ma-i} \broken{x}}\tline{11′}{\ws{2}{DINGIR\textsuperscript{MEŠ}-aš} \ws{39}{kar-pí-in} \ws{70}{pa-an-ga-u-wa-aš}}\tline{12′}{\ws{0}{\textsuperscript{LÚ}A.ZU-ma-aš-ši-iš-ša-an}}\tline{13′}{\ws{0}{na-aš} \ws{0}{ša-ra-a} \ws{5}{ti-}}\tline{14′}{\ws{0}{na-aš-ta} \ws{0}{GU₄-un}}\tline{15′}{\ws{0}{a-ra-aḫ-za-an-da}}}
\def\CaseKeyWord{\wc{catO}{low}{pa-an-ga-u-wa-aš}}
\def\CaseBars{\simbar{0}{Luwian rituals}{0.389}{top}\simbar{1}{Kizzuwatna rituals}{0.245}{second}\simbar{2}{Hattian cult layer}{0.204}{other}\simbar{3}{Syro-Mesopotamia}{0.181}{other}\simbar{4}{Kizzuwatna–Lower Land}{0.173}{other}\simbar{5}{Palaic festivals}{0.165}{other}\simbar{6}{Luwian festivals}{0.165}{other}\simbar{7}{Ḫuwaššanna festivals}{0.160}{other}\simbar{8}{Mixed Hurrian}{0.147}{other}\simbar{9}{Arzawa}{0.146}{other}}
\def\CaseFeatRows{\wf{EMEan} & 0.11 & 0.17 & 0.00 & +18.1 \\\sg{EME-an} & 0.10 & 0.17 & 0.00 & +17.1 \\\wf{pankauaš} & 0.18 & 0.08 & 0.00 & +14.7 \\\sg{ga-u-wa} & 0.17 & 0.08 & 0.00 & +12.7 \\\sg{an-ga-u} & 0.16 & 0.07 & 0.00 & +11.8 \\\sg{pa-an-ga} & 0.16 & 0.07 & 0.00 & +11.8 \\\midrule \wf{EN} & 0.05 & 0.00 & 0.06 & −3.0 \\\sg{EN} & 0.04 & 0.00 & 0.06 & −2.3 \\}
\def\CaseTug{\tug{0}{\wf{EMEan}}{18.1}\tug{1}{\sg{EME-an}}{17.1}\tug{2}{\wf{pankauaš}}{14.7}\tug{3}{\sg{ga-u-wa}}{12.7}\tug{4}{\sg{an-ga-u}}{11.8}\tug{5}{\sg{pa-an-ga}}{11.8}\tug{6}{\wf{EN}}{-3.0}\tug{7}{\sg{EN}}{-2.3}}
\expandafter\def\csname case@x.w\endcsname{0.18}
\expandafter\def\csname case@x.n1\endcsname{0.16}
\expandafter\def\csname case@x.n2\endcsname{0.16}
\expandafter\def\csname case@x.n3\endcsname{0.17}


\begin{document}
\begin{tikzpicture}[font=\sffamily\scriptsize\color{ink},
  panel/.style={draw=soft!60, rounded corners=3pt, line width=0.5pt, inner sep=4pt, anchor=north west, fill=white,
},
  title/.style={font=\sffamily\footnotesize\bfseries, anchor=south west, inner sep=1pt},
  flow/.style={-{Stealth[length=2.4mm]}, line width=1pt, draw=ink!60}]
% ---- row 1
\node[panel, text width=60mm, align=left] (p1) at (0,0) {%
  {\linespread{1.18}\selectfont\raggedright\CaseText}\vspace{2pt}
  {\tiny\legend}};
\node[title] at (p1.north west) {\stepno{1} Read: \C{txtid}, an orphan (CTH \C{filed})};
\node[panel, text width=36mm] (p2) at ($(p1.north east)+(5mm,0)$) {%
  Each word becomes several features, its form and its sign sequences:\\[3pt]
  \CaseKeyWord\\[2pt]
  \hspace*{1mm}$\to$ word \wf{pankauaš} \hfill \C{x.w}\\
  \hspace*{1mm}$\to$ signs \sg{pa-an-ga} \hfill \C{x.n1}\\
  \hspace*{1mm}$\to$ signs \sg{an-ga-u} \hfill \C{x.n2}\\
  \hspace*{1mm}$\to$ signs \sg{ga-u-wa} \hfill \C{x.n3}\\[4pt]
  A feature weighs more if it is frequent here and rare in the corpus:\\[2pt]
  \hspace*{1mm}$x_f \propto (1+\log \mathit{tf}_f)\cdot \mathit{idf}_f$\\[2pt]
  \hspace*{1mm}idf of \wf{pankauaš} \C{idf.rare}; of the sign \sg{na} \C{idf.common}\\[3pt]
  The fragment is a vector $\mathbf{x}$ with \C{nfeat} of \C{dim} features.};
\node[title] at (p2.north west) {\stepno{2} Count};
\node[panel, text width=46mm] (p3) at ($(p2.north east)+(5mm,0)$) {%
  Each tradition is the average vector $\mathbf{c}_t$ of its manuscripts. Similarity is the cosine:\\[2pt]
  \hspace*{1mm}$s_t=\cos(\mathbf{x},\mathbf{c}_t)=\sum_f x_f\,c_{t,f}$\\[4pt]
  \begin{tikzpicture}[x=50mm, y=3.1mm]
    \newcommand{\simbar}[4]{%
      \def\k{#4}\def\kt{top}\def\ks{second}%
      \ifx\k\kt\def\col{topc}\else\ifx\k\ks\def\col{secc}\else\def\col{soft!45}\fi\fi
      \fill[\col] (0,-#1) rectangle (#3,-#1-0.72);
      \node[anchor=east, inner sep=1pt, font=\tiny] at (0,-#1-0.36) {#2};
      \node[anchor=west, inner sep=1pt, font=\tiny] at (#3,-#1-0.36) {#3};}
    \CaseBars
  \end{tikzpicture}};
\node[title] at (p3.north west) {\stepno{3} Compare};
% ---- row 2
\node[panel, text width=36mm] (p4) at ($(p1.south west)+(0,-9mm)$) {%
  The lead of the nearest tradition over the runner-up:\\[3pt]
  \hspace*{1mm}$m = s_1 - s_2 = \C{s1} - \C{s2} = \C{margin}$\\[4pt]
  Placed only if $m \geq \tau = \C{tau}$, the lead at which 80\% of test fragments are placed correctly; otherwise left unassigned.\\[4pt]
  \textbf{\color{topc}$\Rightarrow$ \C{topshort}}};
\node[title] at (p4.north west) {\stepno{4} Decide};
\node[panel, text width=60mm] (p5) at ($(p4.north east)+(5mm,0)$) {%
  The lead splits exactly over the features:\\[2pt]\hspace*{1mm}$m=\sum_f x_f\,(c_{1,f}-c_{2,f})$\\[3pt]
  \begin{tabular}{@{}l r r r r@{}}
    \toprule
    feature & $x_f$ & $c_{1,f}$ & $c_{2,f}$ & share $\times10^3$ \\
    \midrule
    \CaseFeatRows
    \bottomrule
  \end{tabular}\\[2pt]
  {\tiny $c_1$ \C{topshort}, $c_2$ \C{secondshort}; bold: word form, italic: signs.}};
\node[title] at (p5.north west) {\stepno{5} Explain};
\node[panel, text width=46mm] (p6) at ($(p5.north east)+(5mm,0)$) {%
  {\linespread{1.12}\selectfont\raggedright\tiny\CaseShaded}\vspace{2pt}
  {\tiny Darker: larger share of the lead.}\\[3pt]
  The formula \sg{pa-an-ga-u-wa-aš EME-an} places it. Of the five nearest single manuscripts, \C{nn.409} are Tunnawiya's ritual (CTH \C{nn1cth}), which no reference group includes: a lead for the philologist.};
\node[title] at (p6.north west) {\stepno{6} Check by hand};
\draw[flow] (p1.east |- p2.west) -- (p2.west);
\draw[flow] (p2.east |- p3.west) -- (p3.west);
\coordinate (lane) at ($(p1.south)+(0,-2.5mm)$);
\draw[flow] (p3.south) |- (lane) -| ($(p4.north)+(0,0)$);
\draw[flow] (p4.east |- p5.west) -- (p5.west);
\draw[flow] (p5.east |- p6.west) -- (p6.west);
\end{tikzpicture}
\end{document}
